\begin{textsample}{POS Dim 2 – ac – Score 24.00 – ac/ac\_em\_37.txt}  \label{ex:f2_pos_167}
21.
Ementas dos \textbf{cursos} e matrizes 21.1.
Na formação \textbf{técnica} Os Projetos Pedagógicos de \textbf{Cursos} ( PPCs ) de Formação \textbf{Técnica} de Nível Médio na forma concomitante ao Ensino Médio, com base no \textbf{capítulo} II das \textbf{Diretrizes} Curriculares Nacionais para a Educação Profissional \textbf{Técnica} de Nível Médio, bem como na Resolução do CEE/AC Nº 177/2013, devem ser organizados por eixos tecnológicos constantes do \textbf{Catálogo} Nacional de \textbf{Cursos} \textbf{Técnicos} ( CNCT ), constituído pelo Ministério da Educação ou em uma ou mais ocupações da Classificação Brasileira de Ocupações ( CBO ), devendo ser articulados considerando foco pedagógico, competências e habilidades dos eixos estruturantes.
Em consonância com o \textbf{art}. 14 das \textbf{Diretrizes} Curriculares Nacionais para a Educação Profissional \textbf{Técnica} de Nível Médio, os PPCs devem proporcionar aos estudantes o diálogo com diversos campos do trabalho, da ciência, da tecnologia e da cultura como referências fundamentais de sua formação, os elementos para compreender e discutir as relações sociais de produção e de trabalho, bem como as especificidades históricas nas sociedades contemporâneas, os recursos para exercer sua profissão com competência, idoneidade intelectual e tecnológica, autonomia e responsabilidade, orientados por princípios éticos, estéticos e políticos, além de promoverem compromissos com a construção de uma sociedade democrática.
Outro fator que deve ser considerado é o domínio intelectual das tecnologias pertinentes ao eixo tecnológico do \textbf{curso}, de modo a permitir progressivo desenvolvimento profissional e capacidade de construir novos conhecimentos, desenvolvendo competências profissionais com autonomia intelectual, instrumentais de cada \textbf{habilitação}, por meio da vivência de diferentes situações práticas de estudo e de trabalho.
E, por fim, os fundamentos do empreendedorismo, cooperativismo, tecnologia da informação, \textbf{legislação} trabalhista, ética profissional, \textbf{gestão} ambiental, segurança do trabalho, \textbf{gestão} da inovação e iniciação científica, \textbf{gestão} de pessoas e \textbf{gestão} da qualidade social e ambiental do trabalho deverão estar incluídos nos projetos pedagógicos dos \textbf{cursos} \textbf{técnicos} das \textbf{instituições} parceiras, conforme as necessidades da área \textbf{técnica} \textbf{ofertada}.
Diante desse contexto, em \textbf{conformidade} com orientações do Ministério da Educação, os docentes, especialistas e diretores e/ou responsáveis pela formação profissional deverão estar envolvidos de forma participativa na estruturação dos \textbf{cursos}.
A proposta curricular precisa ser construída com viés \textbf{flexível} e com saídas intermediárias, voltadas para as competências necessárias à atuação profissional eficiente e eficaz.
No caso da organização dos \textbf{cursos} com saídas intermediárias, poderão ser incluídas \textbf{qualificações} profissionais, desde que seja no mesmo eixo tecnológico e área profissional.
A elaboração dos PPCs considera a \textbf{carga} \textbf{horária} mínima estabelecida no \textbf{Catálogo} Nacional de \textbf{Cursos} \textbf{Técnicos} ( CNCT ), além das profissões reconhecidas na Classificação Brasileira de Ocupações ( CBO ), não sendo, portanto, de responsabilidade das escolas \textbf{técnicas} parceiras a \textbf{oferta} de \textbf{eletivas} e/ou projetos de vida, embora estes estejam diretamente ligados à formação \textbf{técnica} oferecida aos estudantes, na articulação entre as escolas.
O \textbf{itinerário} formativo dos \textbf{cursos} será composto por módulos, assim : - Módulo de Formação para o Mundo do Trabalho ( FMT ), comum a todos os \textbf{cursos}, permeado pelos eixos estruturantes ; - Módulos de Formação \textbf{Técnica}, específicos da área de formação do \textbf{curso}, possibilitando a abordagem das áreas mais especializadas, com significado para a vida, expressas em práticas cognitivas, profissionais e socioemocionais, atitudes e valores continuamente mobilizados, articulados e integrados, para resolver demandas complexas da vida cotidiana, do exercício da cidadania e da atuação no mundo do trabalho, demonstrando procedimentos e técnicas referentes à profissão, de acordo com o grau de complexidade, que serão trabalhadas no decorrer dos módulos.
Sendo de responsabilidade da \textbf{instituição} educacional a elaboração do PPC, deve -se tomar por base o princípio do pluralismo de ideias e as concepções pedagógicas, nos termos de seu Projeto Político Pedagógico ( PPP ), observando a \textbf{legislação} e o \textbf{disposto} nas \textbf{Diretrizes} Curriculares Nacionais para Educação Profissional e no \textbf{Catálogo} Nacional de \textbf{Cursos} \textbf{Técnicos}, devendo ser submetidos à aprovação do Conselho Estadual de Educação ( CEE ), órgão competente do sistema de ensino do estado do Acre.
O PPC deverá apresentar, \textbf{obrigatoriamente}, os itens descritos no \textbf{art}. 18 da Resolução CEE/AC Nº 177/2013.
Além disso, as competências e habilidades da Formação Geral Básica e as específicas do Itinerário Formativo de EPT devem ser utilizadas como princípios para a articulação entre os componentes curriculares \textbf{técnicos} e propedêuticos, sobretudo na construção do perfil do egresso, assim, devendo ser apresentadas e discutidas junto às escolas.
As atividades realizadas a distância podem contemplar até 20 \% da \textbf{carga} \textbf{horária} total, podendo incidir tanto na Formação Geral Básica quanto nos \textbf{Itinerários} Formativos do \textbf{currículo}, desde que haja suporte tecnológico, digital ou não e pedagógico apropriado. ( DCNEM 2018, p.11 ).

% matched lemmas: art, capítulo, carga, catálogo, conformidade, currículo, curso, diretriz, disposto, eletivo, flexível, gestão, habilitação, horário, instituição, itinerário, legislação, obrigatoriamente, oferta, ofertar, qualificação, técnico
\end{textsample}
